\chapter{Product Design Specification}
\label{ch:specification}

\begin{guidance}[What this chapter does]
The specification is the control document for the whole design: it defines
what the product must do, within which constraints, and how success will be
tested. It should not predict the design. Use short, definitive statements
and quantify every parameter you can. Show how the specification was
justified (literature, market research, user survey) and how it was reviewed
and updated since the Progress Report. Typical length: 400--700 words.

Number every requirement with \verb|\req{req:label}| in the first column. The
IDs R1, R2 \ldots\ are generated automatically and can be referenced
anywhere with \verb|\cref{req:label}|, for example in the test table of
\cref{ch:implementation} and the evaluation in \cref{ch:evaluation}. This
traceability from requirement to test to evaluation is exactly what the
markers look for.

Further reading: \textcite{cross2008}, \textcite{ulrich2012}; on writing
testable requirements: \textcite{dick2017}.
\end{guidance}

\section{Sources of the requirements}
\label{sec:spec-sources}

\section{Requirements}
\label{sec:requirements}

\Cref{tab:requirements} lists the requirements. Requirements marked
\enquote{D} (demand) must be met; \enquote{W} (wish) are desirable.

\begin{table}[htbp]
  \centering
  \caption{Product design specification (example)}
  \label{tab:requirements}
  \small
  \begin{tabular}{p{0.06\linewidth}p{0.14\linewidth}p{0.33\linewidth}cp{0.2\linewidth}}
    \toprule
    ID & Category & Requirement & D/W & Verification \\
    \midrule
    \req{req:accuracy} & Performance & Temperature measured with an accuracy
      of $\pm\SI{0.5}{\celsius}$ between \SIrange{0}{40}{\celsius}
      & D & Comparison with a calibrated reference \\
    \req{req:interval} & Performance & Logging interval adjustable from
      \SIrange{1}{60}{\minute} & D & Functional test \\
    \req{req:battery} & Power & Battery life of at least 30 days at a
      \SI{5}{\minute} interval & D & Measured current, calculation \\
    \req{req:cost} & Cost & Component cost below £40 per unit
      & D & Bill of materials \\
    \req{req:size} & Size & Fits within a $\SI{80}{\milli\metre} \times
      \SI{50}{\milli\metre} \times \SI{25}{\milli\metre}$ enclosure
      & W & Measurement \\
    \bottomrule
  \end{tabular}
\end{table}

\section{Changes since the Progress Report}
\label{sec:spec-changes}

\begin{guidance}
Which requirements changed, why, and when (refer to the logbook entry).
For example: \enquote{\Cref{req:battery} was relaxed from 60 to 30 days after
the current measurements in week 18 (logbook, p.~42).}
\end{guidance}
